\scanpage{35}
l'on doit avoir, si $a\in V\cap T'(W)$, $R_K\subset G_{x\,K}$. L'inverse est vrai, car comme $T_K$ est fid. plat par $K$, on voit par descente que les pts (à coeff. ds des $K$-alg. quelconques $\Lambda$) $s.x$ $(s\in T_K(\Lambda))$ proviennent d'un pt bien déterminé de $V_K$. Or soit \deleted{\unsure{$\xi$}}
\[
S_K=\overline{\varphi(R_K)}\qquad p\colon G_K\unsure{\mathrel{-}|\longrightarrow}G,
\]
et soit $H$ le ss-gpe algébrique \inserted{fermé} de $G$ engendré par $S$. On a donc
\[
x\in V\cap T'(V)\quad\Longleftrightarrow\quad x\in V_{\deleted{\unsure{x|0}}}^{H}\quad \deleted{\unsure{H|\forall}\unsure{\subset|\in}\unsure{G}\ \unsure{H_x|H}\ \text{\unsure{tous|tels}}}.
\]

Par suite, pour que \deleted{les invariants} pour tout $V$, on ait $V\cap T'(V)=V^G$, i.e. que le foncteur $\mathrm{Rep}(G)=\mathcal M\longrightarrow$ (structures vectorielles mixtes pour $L(K/k)$ soit pl. fid, il f. et s. que pour $\forall\,V$, \deleted{les invariants} on ait
\[
V^H=V^G.
\]

Supposons $P$ quelconque. On a l'hom canonique de bitorseurs
\[
\deleted{\unsure{\underline{\mathrm{Isom}}|{\mathrm{Isom}}}_{\unsure{\otimes}}(\unsure{T'|T},T_K))}
\]
\diagram{\begin{tikzcd}[ampersand replacement=\&,row sep=small]
P\otimes_K P\arrow[r] \& G_K\\
(v,u)\arrow[r,mapsto] \& v^{-1}u
\end{tikzcd}}
\deleted{induisant}
\[
\deleted{P_L^G\times P_L\mathrel{-}G_L}
\]
donc $\varphi\colon\mathrm{Spec}\,L\longrightarrow P$ (un $\mathrm{Spec}\,K$-m., bien sûr) s'en déduit \deleted{de, grâce à l'isom $P_L\simeq G_L$}
\[
\mathrm{Spec}(L)\mathbin{\unsure{\otimes|\times}}_{\mathrm{Spec}\,K}\mathrm{Spec}(L)\mathrel{\unsure{-|\longrightarrow}}
P\mathbin{\unsure{\otimes|\times}}_{\mathrm{Spec}\,\unsure{K|k}}P\longrightarrow G_K\,;
\]
l'image fermée $\deleted{\unsure{P|R}}R_K$ est donc définie intrinsèquement. On vérifie encore (par descente \unsure{platisation|fidèlement|plate}) que \deleted{\unsure{p}} l'on a, pour tt $G$-module $V$
\[
\deleted{x\in}V\cap T'(V)=\{\,x\in V,\mid R_K\subset G_{x\,K}\,\}.
\]
Définissant $H$ comme ci-dessus, on trouve donc
\[
V\cap T'(W)=V^H,
\]
et la $\overline{\mathrm m}$ conclusion que plus haut ....
% paper watermark, sideways
{\centering \textsc{\unsure{Extra Strong}}\\
\textit{\unsure{Ma}\unsure{in|nn}\unsure{ol|on}\unsure{uy|ny|y} \unsure{1808}}\par}
% archive sheet number: 34
